%% 1-20 角の 大きさ（10ページ、p.202〜211）
%% 第1巻のきまり：15 pt・行送り28 pt、分かち書き（文節の間に半角の空白）、漢字はすべてルビ、式は全角の＋＝で書く
%% 小4の内容を、小3向けに早めにやさしく扱う。直角・辺・ちょう点・三角じょうぎは 1-19 と同じ ことばと絵にそろえる。

\makeatletter
% 直角のしるし（1-19 の \bkqrright と同じ絵）  \bkqsright{(かど)}{1つめの辺の向き（度）}{大きさ}
\newcommand\bkqsright[3]{\begin{scope}[shift={#1}, rotate=#2]
  \draw[bkTheme, line width=0.9pt] (#3,0) -- (#3,#3) -- (0,#3);\end{scope}}
% 三角じょうぎ（1-19 の \bkqrjougi と同じ絵）  \bkqsjougi{(直角のかど)}{1:45°の形 / 2:30°と60°の形}{横の辺の長さ}{回す角度}
\newcommand\bkqsjougi[4]{\begin{scope}[shift={#1}, rotate=#4]
  \pgfmathsetmacro\bkqs@a{#3}%
  \ifnum#2=1 \pgfmathsetmacro\bkqs@b{#3}\else\pgfmathsetmacro\bkqs@b{#3*0.5774}\fi
  \pgfmathsetmacro\bkqs@r{(\bkqs@a+\bkqs@b-sqrt(\bkqs@a*\bkqs@a+\bkqs@b*\bkqs@b))/2}%
  \pgfmathsetmacro\bkqs@k{0.42}%
  \fill[bkBlue!14, even odd rule] (0,0) -- (\bkqs@a,0) -- (0,\bkqs@b) -- cycle
    ({\bkqs@r*(1-\bkqs@k)},{\bkqs@r*(1-\bkqs@k)}) -- ({\bkqs@r+(\bkqs@a-\bkqs@r)*\bkqs@k},{\bkqs@r*(1-\bkqs@k)})
    -- ({\bkqs@r*(1-\bkqs@k)},{\bkqs@r+(\bkqs@b-\bkqs@r)*\bkqs@k}) -- cycle;
  \draw[bkBlue!70!black, line width=0.9pt, line join=round] (0,0) -- (\bkqs@a,0) -- (0,\bkqs@b) -- cycle;
  \draw[bkBlue!55!black, line width=0.6pt, line join=round]
    ({\bkqs@r*(1-\bkqs@k)},{\bkqs@r*(1-\bkqs@k)}) -- ({\bkqs@r+(\bkqs@a-\bkqs@r)*\bkqs@k},{\bkqs@r*(1-\bkqs@k)})
    -- ({\bkqs@r*(1-\bkqs@k)},{\bkqs@r+(\bkqs@b-\bkqs@r)*\bkqs@k}) -- cycle;
  \pgfmathtruncatemacro\bkqs@n{(\bkqs@a-6)/2}%
  \foreach \i in {1,...,\bkqs@n} {\pgfmathtruncatemacro\bkqs@t{mod(\i,5)==0}%
    \ifnum\bkqs@t=1 \draw[bkBlue!60!black, line width=0.4pt] (2*\i,0) -- (2*\i,1.6);
    \else\draw[bkBlue!60!black, line width=0.3pt] (2*\i,0) -- (2*\i,0.9);\fi}
\end{scope}}
% 形や角につける名前の札（1-19 と同じ）  \bkqstag{(x,y)}{A}
\newcommand\bkqstag[2]{\node[circle, draw=bkInk!70, fill=white, line width=0.7pt, text=bkInk,
  font=\sffamily\bfseries\footnotesize, minimum size=5.2mm, inner sep=0pt] at #1 {#2};}
% 分度器（\bkprotractor と同じ色・目もりで、数字を2れつ入れたもの）
%   \bkqsbundoki[外れつの字の大きさ pt]{(中心)}{半径}{数字を書く角度のリスト}{はかる辺の向きのリスト}
%   外がわの数字は右の0から、内がわの数字は左の0から かぞえる。辺は数字の下を通す（数字に地の色のふち）。
%   0 と 180 の数字は、0の線に かくれないよう 線の外がわ（下の おびの中）に よこ書きで おく。
\newcommand\bkqsbundoki[5][6.5]{\begin{scope}[shift={#2}]
  \pgfmathsetmacro\bkqs@sb{#3*0.09+0.6}%
  \fill[bkBlue!8] (-#3,0) arc (180:0:#3) -- (#3,-\bkqs@sb) -- (-#3,-\bkqs@sb) -- cycle;
  \draw[bkBlue!70!black, line width=0.8pt] (-#3,0) arc (180:0:#3) -- (#3,-\bkqs@sb) -- (-#3,-\bkqs@sb) -- cycle;
  \draw[bkBlue!50!black, line width=0.4pt] (-#3,0) -- (#3,0);
  \draw[bkBlue!50!black, line width=0.5pt] (0,0) ++(-#3*0.45,0) arc (180:0:#3*0.45);
  \foreach \d in {0,...,180} {\pgfmathtruncatemacro\bkqs@t{mod(\d,10)==0}\pgfmathtruncatemacro\bkqs@f{mod(\d,5)==0}
    \ifnum\bkqs@t=1 \draw[bkInk, line width=0.5pt] (\d:#3) -- (\d:#3*0.86);
    \else\ifnum\bkqs@f=1 \draw[bkInk, line width=0.35pt] (\d:#3) -- (\d:#3*0.91);
    \else\draw[bkInk!70, line width=0.25pt] (\d:#3) -- (\d:#3*0.95);\fi\fi}
  \foreach \a in {#5} {\draw[bkInk!85, line width=1.4pt, line cap=round] (0,0) -- (\a:#3*1.2);}
  \pgfmathsetmacro\bkqs@fo{#1}\pgfmathsetmacro\bkqs@fi{#1*0.85}%
  \pgfmathsetmacro\bkqs@ro{#3*0.86-0.6-\bkqs@fo*0.13}\pgfmathsetmacro\bkqs@ri{\bkqs@ro-1.7-\bkqs@fo*0.24}%
  \xdef\bkqsro{\bkqs@ro}%
  \foreach \d in {#4} {\pgfmathtruncatemacro\bkqs@in{180-\d}
    \ifnum\d=0
      \node[font=\sffamily\bfseries\fontsize{\bkqs@fo}{\bkqs@fo}\selectfont, text=bkInk, fill=bkBlue!8, inner sep=0.3pt] at (\bkqs@ro,-\bkqs@sb/2) {\d};
      \node[font=\sffamily\fontsize{\bkqs@fi}{\bkqs@fi}\selectfont, text=bkBlue!65!black, fill=bkBlue!8, inner sep=0.3pt] at (\bkqs@ri,-\bkqs@sb/2) {\bkqs@in};
    \else\ifnum\d=180
      \node[font=\sffamily\bfseries\fontsize{\bkqs@fo}{\bkqs@fo}\selectfont, text=bkInk, fill=bkBlue!8, inner sep=0.3pt] at (-\bkqs@ro,-\bkqs@sb/2) {\d};
      \node[font=\sffamily\fontsize{\bkqs@fi}{\bkqs@fi}\selectfont, text=bkBlue!65!black, fill=bkBlue!8, inner sep=0.3pt] at (-\bkqs@ri,-\bkqs@sb/2) {\bkqs@in};
    \else
      \node[font=\sffamily\bfseries\fontsize{\bkqs@fo}{\bkqs@fo}\selectfont, text=bkInk, fill=bkBlue!8, inner sep=0.3pt, rotate=\d-90] at (\d:\bkqs@ro) {\d};
      \node[font=\sffamily\fontsize{\bkqs@fi}{\bkqs@fi}\selectfont, text=bkBlue!65!black, fill=bkBlue!8, inner sep=0.3pt, rotate=\d-90] at (\d:\bkqs@ri) {\bkqs@in};
    \fi\fi}
  \draw[bkInk, line width=0.5pt] (0,-1.6) -- (0,1.6);
  \fill[bkRed] (0,0) circle (0.7);\end{scope}}
% 角（2本の辺と、まわった大きさを示す弧）  \bkqskaku{(ちょう点)}{1本めの向き}{2本めの向き}{辺の長さ}{弧の半径}
\newcommand\bkqskaku[5]{\begin{scope}[shift={#1}]
  \draw[bkInk!85, line width=1.3pt, line cap=round] (#2:#4) -- (0,0) -- (#3:#4);
  \draw[bkTheme, line width=1pt] (#2:#5) arc (#2:#3:#5);
  \fill[bkInk] (0,0) circle (0.8);\end{scope}}
% とびら（上から見た図）  \bkqsdoor{(x,y)}{ひらいた角度}
\newcommand\bkqsdoor[2]{\begin{scope}[shift={#1}]
  \fill[white, rounded corners=2mm] (0,-3) rectangle (40,24);
  \draw[bkGray, line width=2.6pt] (2,1) -- (14,1) (32,1) -- (38,1);
  \draw[bkGray!70, line width=0.6pt, dash pattern=on 1.5pt off 1.5pt] (14,1) -- (32,1);
  \draw[bkTheme, line width=1.1pt, -{Stealth[length=4pt]}] (14,1) ++(0:12) arc (0:#2-6:12);
  \draw[bkStem, line width=2.8pt, line cap=round] (14,1) -- ++(#2:18);
  \fill[bkInk] (14,1) circle (1.1);\end{scope}}
% 上から見た子ども（頭とかた。向いている方へ はなが出る）  \bkqstop{(x,y)}{向き（度）}
\newcommand\bkqstop[2]{\begin{scope}[shift={#1}, rotate=#2-90]
  \fill[bkRoute6] (0,-0.6) ellipse (6.5 and 3);
  \fill[bkHair] (0,0) circle (3.6);
  \fill[bkSkin] (0,3.5) ellipse (1.3 and 1);\end{scope}}
% 1-19 で おって作った 直角の紙  \bkqspaper{(直角のかど)}{回す角度}{色}
\newcommand\bkqspaper[3]{\begin{scope}[shift={#1}, rotate=#2]\colorlet{bkqs@pc}{#3}%
  \filldraw[fill=bkqs@pc, draw=bkqs@pc!60!black, line width=0.7pt, line join=round] (0,0) -- (14,0) -- (12.5,8) -- (7,13.5) -- (0,13.5) -- cycle;
  \end{scope}}
\makeatother

\bkchapter{1}{20}{\R{角}{かく}の \R{大}{おお}きさ}{\bklevel{1}}{かくの おおきさ}
\bkmichishirube{1-10,1-12,1-19}{4-2,4-8,5-15}

\begin{goal}
\begin{itemize}[label={\color{bkTheme}●}, leftmargin=1.4\zw, itemsep=0pt, topsep=0pt]
\item \R{角}{かく}の \R{大}{おお}きさは 「まわった \R{大}{おお}きさ」と わかる。
\item 1\R{回転}{かい|てん}は 360°、 \R{直角}{ちょっ|かく}は 90°と わかる。
\item \R{分度器}{ぶん|ど|き}で \R{角}{かく}の \R{大}{おお}きさを はかれる。
\end{itemize}
\end{goal}

\begin{hatena}
\begin{bkfigure}
\begin{tikzpicture}[x=1mm, y=1mm]
  % A：その ばで 1回転
  \fill[white, rounded corners=3mm] (0,0) rectangle (59,56);
  \node[circle, fill=bkYellow, text=white, font=\sffamily\bfseries, minimum size=6.5mm, inner sep=0pt] at (6,50) {A};
  \node[anchor=west, font=\small, text=bkInk] at (11,50) {その ばで くるっと};
  \draw[bkYellow!80!black, line width=1.2pt, dash pattern=on 2pt off 1.5pt] (49,11) arc (0:180:19.5 and 6);
  \bkkid{(29.5,7)}{36}{bkRoute6}{bkHair}
  \draw[bkYellow!80!black, line width=1.6pt, -{Stealth[length=7pt]}] (10,11) arc (180:350:19.5 and 6);
  % B：長い はりが 1回転
  \begin{scope}[shift={(63,0)}]
  \fill[white, rounded corners=3mm] (0,0) rectangle (59,56);
  \node[circle, fill=bkYellow, text=white, font=\sffamily\bfseries, minimum size=6.5mm, inner sep=0pt] at (6,50) {B};
  \node[anchor=west, font=\small, text=bkInk] at (11,50) {\R{時計}{と|けい}の \R{長}{なが}い はり};
  \bkclock{(29.5,22)}{15}{10}{8}
  \draw[bkYellow!80!black, line width=1.6pt, -{Stealth[length=7pt]}] (29.5,22) ++(36:19) arc (36:-309:19);
  \end{scope}
\end{tikzpicture}
\end{bkfigure}
\noindent A も B も、 ぐるっと 1\R{回転}{かい|てん}して もとに もどります。\\
まわった \R{大}{おお}きさも、 \R{数}{かず}で あらわせるかな？\\
1\R{回転}{かい|てん}すると、 どれだけ まわった ことに なるのかな？
\end{hatena}

\clearpage
%% ---------------- p.203 やってみよう ----------------
\begin{yattemiyou}[からだで まわって みよう]
\begin{enumerate}[label={\bkcircled[bkGreen]{\arabic*}}, leftmargin=2.2\zw, itemsep=2pt, topsep=0pt]
\item こくばんを \R{向}{む}いて \R{立}{た}とう。
\item その ばで、 \R{右}{みぎ}を \R{向}{む}こう。 かどの \R{形}{かたち}（\R{直角}{ちょっ|かく}）の ぶんだけ まわったね。
\item また \R{右}{みぎ}を \R{向}{む}こう。 \R{何回}{なん|かい}で こくばんの \R{向}{む}きに もどるかな？
\end{enumerate}
\begin{bkfigure}
\begin{tikzpicture}[x=1mm, y=1mm]
  \fill[white, rounded corners=3mm] (0,0) rectangle (122,69);
  % へや（上から見た図）
  \draw[bkGray!80, line width=1pt] (21,6) rectangle (101,56);
  \fill[bkGreen!60!black] (46,53.5) rectangle (76,56);
  \node[anchor=south, font=\sffamily\bfseries\small] at (61,56.5) {こくばん};
  \fill[bkBlue!35] (98.5,16) rectangle (101,46);
  \node[anchor=west, font=\sffamily\bfseries\small] at (102,31) {まど};
  \fill[bkStem!60] (46,6) rectangle (76,8.5);
  \node[anchor=north, font=\sffamily\bfseries\small] at (61,5.5) {たな};
  \fill[bkYellow!50] (21,20) rectangle (23.5,42);
  \node[anchor=east, font=\sffamily\bfseries\small] at (20,31) {ろうか};
  % 子ども
  \bkqstop{(61,31)}{90}
  \draw[bkTheme, line width=1.4pt, -{Stealth[length=6pt]}] (61,31) ++(85:13) arc (85:5:13);
  \node[anchor=south west, font=\sffamily\bfseries\footnotesize, text=bkTheme] at (74,40) {\R{直角}{ちょっ|かく} 1つ\R{分}{ぶん}};
  \node[anchor=north east, font=\sffamily\footnotesize, text=bkGray] at (120,67.5) {\R{上}{うえ}から \R{見}{み}た \R{図}{ず}};
\end{tikzpicture}
\end{bkfigure}
\noindent\hspace*{1\zw}\R{右}{みぎ}を\enspace\bkkakikomi\enspace \R{回}{かい} \R{向}{む}くと、 もとの \R{向}{む}きに もどる。\par
\noindent\hspace*{1\zw}{\small\color{bkOrange!80!black}めが まわらないように、 ゆっくり まわろう。}
\end{yattemiyou}

\bkfukidashi{sora}{4\R{回}{かい}で こくばんの \R{向}{む}きに もどったよ！\\\R{直角}{ちょっ|かく} 4つ\R{分}{ぶん}で、 ぐるっと 1\R{回転}{かい|てん}だね。}

\clearpage
%% ---------------- p.204〜205 数学の道具 ----------------
\bksection{doguu}{\R{数学}{すう|がく}の \R{道具}{どう|ぐ}}

\bksubsection{1}{\R{角}{かく}は 「まわった \R{大}{おお}きさ」}

\label{bk:1-20:kaku}とびらを あけて みよう。 とびらは、 ちょうつがいを \R{中心}{ちゅう|しん}に まわります。

\begin{bkfigure}
\begin{tikzpicture}[x=1mm, y=1mm]
  \bkqsdoor{(0,0)}{30}\node[anchor=north, font=\sffamily\bfseries\footnotesize] at (20,-3.5) {すこし あける};
  \bkqsdoor{(45,0)}{90}\bkqsright{(59,1)}{0}{3}\node[anchor=north, font=\sffamily\bfseries\footnotesize] at (65,-3.5) {\R{直角}{ちょっ|かく}まで あける};
  \bkqsdoor{(90,0)}{135}\node[anchor=north, font=\sffamily\bfseries\footnotesize] at (110,-3.5) {\R{大}{おお}きく あける};
\end{tikzpicture}
\end{bkfigure}

\noindent\begin{minipage}[c]{0.64\linewidth}
とびらの ちょうつがいの ように、 1つの \R{点}{てん}から 2\R{本}{ほん}の \R{線}{せん}が \R{出}{で}て いる \R{形}{かたち}を \R{角}{かく}と いいます。\bkidx{かく}{\R{角}{かく}} どれだけ ひらいたか（まわったか）が、 \R{角}{かく}の \R{大}{おお}きさです。
\end{minipage}\hfill
\begin{minipage}[c]{0.34\linewidth}\centering
\begin{tikzpicture}[x=1mm, y=1mm]
  \bkqskaku{(0,0)}{0}{50}{30}{9}
  \node[font=\sffamily\bfseries\footnotesize, text=bkTheme] at (25:13.5) {\R{角}{かく}};
  \node[anchor=north, font=\sffamily\bfseries\footnotesize, text=bkInk] at (0,-1.5) {ちょう\R{点}{てん}};
  \node[anchor=north, font=\sffamily\bfseries\footnotesize, text=bkInk] at (24,-0.5) {\R{辺}{へん}};
  \node[anchor=east, font=\sffamily\bfseries\footnotesize, text=bkInk] at (50:25) {\R{辺}{へん}};
\end{tikzpicture}
\end{minipage}

\bksubsection{2}{\R{直角}{ちょっ|かく} 4つで 1\R{回転}{かい|てん}}

\label{bk:1-20:kaiten}\R{直角}{ちょっ|かく}の \R{紙}{かみ}（1-19）を、 かどを \R{合}{あ}わせて ならべよう。

\begin{bkfigure}
\begin{tikzpicture}[x=1mm, y=1mm]
  \bkqspaper{(22,16)}{0}{bkRoute6!35}\bkqspaper{(22,16)}{90}{bkRoute2!35}
  \bkqspaper{(22,16)}{180}{bkRoute4!35}\bkqspaper{(22,16)}{270}{bkRoute5!30}
  \fill[bkRed] (22,16) circle (1);
  \node[anchor=west, font=\sffamily\bfseries\small, align=left] at (40,16) {\R{直角}{ちょっ|かく} 4つで\\1\R{回転}{かい|てん}};
  \bkqspaper{(94,2)}{0}{bkRoute6!35}\bkqspaper{(94,2)}{90}{bkRoute2!35}
  \draw[bkInk!70, line width=0.8pt] (76,2) -- (112,2);
  \fill[bkRed] (94,2) circle (1);
  \node[anchor=north, font=\sffamily\bfseries\small] at (94,0) {\R{直角}{ちょっ|かく} 2つで \R{半回転}{はん|かい|てん}（まっすぐ）};
\end{tikzpicture}
\end{bkfigure}

\noindent \R{直角}{ちょっ|かく} 4つ\R{分}{ぶん}で、 ぐるっと 1\R{回転}{かい|てん}です。\bkidx{かいてん}{\R{回転}{かい|てん}}

\clearpage
\bksubsection{3}{\R{角}{かく}の \R{大}{おお}きさの たんい「\R{度}{ど}」}

\label{bk:1-20:do}1\R{回転}{かい|てん}を 360に \R{分}{わ}けた 1つ\R{分}{ぶん}を 1\R{度}{ど}と いい、 1° と かきます。\bkidx{ど}{\R{度}{ど}（°）}

\bkkigocards{\bkkigocard{\tikz[baseline=-0.55em]\draw[line width=0.09em] (0,0) circle (0.17em);}{ど}{\R{角}{かく}の \R{大}{おお}きさの たんい}{degree}{1-20}}{%
\begin{tikzpicture}[x=1mm, y=1mm]
  \bkqskaku{(4,6)}{0}{90}{12}{5}\bkqsright{(4,6)}{0}{2.6}
  \node[font=\sffamily\bfseries\small, text=bkTheme] at (12,14.5) {90°};
  \node[anchor=north, font=\sffamily\bfseries\footnotesize] at (8,0.4) {\R{直角}{ちょっ|かく}};
  \bkqskaku{(32,6)}{0}{180}{9}{5}
  \node[font=\sffamily\bfseries\small, text=bkTheme] at (32,14.5) {180°};
  \node[anchor=north, font=\sffamily\bfseries\footnotesize] at (32,0.4) {\R{半回転}{はん|かい|てん}};
  \draw[bkInk!85, line width=1.3pt, line cap=round] (52,6) -- (61,6);\fill[bkInk] (52,6) circle (0.8);
  \draw[bkTheme, line width=1pt] (52,6) circle (5);
  \node[font=\sffamily\bfseries\small, text=bkTheme] at (52,14.5) {360°};
  \node[anchor=north, font=\sffamily\bfseries\footnotesize] at (52,0.4) {1\R{回転}{かい|てん}};
\end{tikzpicture}}

\noindent 360こは こまかすぎて \R{絵}{え}に かけないので、 10こずつ まとめた 「10°の かけら」で \R{見}{み}て みよう。

\noindent\begin{minipage}[c]{0.4\linewidth}\centering
\begin{tikzpicture}[x=1mm, y=1mm]
  \foreach \k in {0,...,35} {\pgfmathsetmacro\aa{\k*10}
    \ifnum\k<9 \filldraw[fill=bkTheme!45, draw=white, line width=0.5pt] (0,0) -- (\aa:17) arc (\aa:\aa+10:17) -- cycle;
    \else \filldraw[fill=bkBlue!14, draw=white, line width=0.5pt] (0,0) -- (\aa:17) arc (\aa:\aa+10:17) -- cycle;\fi}
  \draw[bkInk!70, line width=0.8pt] (0,0) circle (17);
  \bkqsright{(0,0)}{0}{3}
  \fill[bkInk] (0,0) circle (0.8);
  \node[anchor=south west, font=\sffamily\bfseries\fontsize{8}{9}\selectfont, text=bkTheme!70!black, inner sep=0.8pt] at (40:17.8) {9こ ＝ 90°};
\end{tikzpicture}
\end{minipage}\hfill
\begin{minipage}[c]{0.58\linewidth}
\begin{enumerate}[label={\bkcircled[bkTheme]{\arabic*}}, leftmargin=2\zw, itemsep=0pt, topsep=0pt]
\item 1\R{回転}{かい|てん}は、 10°の かけらが 36こ（10が 36こで 360）。
\item \R{直角}{ちょっ|かく}は、 1\R{回転}{かい|てん}を 4つに \R{同}{おな}じに \R{分}{わ}けた 1つ（\R{直角}{ちょっ|かく} 4つで 1\R{回転}{かい|てん}）。
\item 36を 4つに \R{同}{おな}じに \R{分}{わ}けると 9（\mbox{9 × 4 ＝ 36}）。 \R{直角}{ちょっ|かく}は かけら 9こで、 10が 9こで 90°。
\end{enumerate}
\end{minipage}

\par\vspace{4pt}
\noindent \bkhatenaref の \R{答}{こた}え：1\R{回転}{かい|てん}は 360°。 \R{直角}{ちょっ|かく}は 90°で、\\
\mbox{90 ＋ 90 ＋ 90 ＋ 90 ＝ 360} です（10が \mbox{9 × 4 ＝ 36}こ）。

\bksubsection{4}{\R{分度器}{ぶん|ど|き}で はかる}

\label{bk:1-20:bundoki}\noindent\begin{minipage}[c]{0.5\linewidth}\centering
\begin{tikzpicture}[x=1mm, y=1mm]
  \bkqsbundoki[5.4]{(0,0)}{30}{0,10,...,180}{0,50}
  \draw[bkRed, line width=1pt] (50:\bkqsro) circle (2.4);
  \node[anchor=north, font=\sffamily\bfseries\footnotesize, text=bkRed!80!black] at (0,-4.6) {\R{中心}{ちゅう|しん}};
  \node[anchor=north, font=\sffamily\bfseries\footnotesize, text=bkInk] at (22,-4.6) {0の \R{線}{せん}};
\end{tikzpicture}
\end{minipage}\hfill
\begin{minipage}[c]{0.48\linewidth}
\begin{enumerate}[label={\bkcircled[bkTheme]{\arabic*}}, leftmargin=2\zw, itemsep=0pt, topsep=0pt]
\item \R{中心}{ちゅう|しん}を ちょう\R{点}{てん}に \R{合}{あ}わせる。
\item 0の \R{線}{せん}を \R{辺}{へん}に \R{合}{あ}わせる。
\item 0から かぞえた \R{目}{め}もりを よむ。 → 50°
\end{enumerate}
\end{minipage}\bkidx{ぶんどき}{\R{分度器}{ぶん|ど|き}}

\par\vspace{4pt}
\noindent \R{分度器}{ぶん|ど|き}には、 \R{数}{かず}が 2れつ あります。
\R{外}{そと}がわの \R{黒}{くろ}い \R{数}{かず}は \R{右}{みぎ}の 0から、 \R{内}{うち}がわの \R{青}{あお}い \R{数}{かず}は \R{左}{ひだり}の 0から ふえて いきます。
\R{辺}{へん}が \R{右}{みぎ}の 0に \R{合}{あ}って いたら \R{外}{そと}がわ、 \R{左}{ひだり}の 0に \R{合}{あ}って いたら \R{内}{うち}がわを よみます。
どちらの \R{辺}{へん}からでも はかれるように、 2れつ あるのです。

\par\vspace{3pt}
\noindent よむ \R{前}{まえ}に、 「\R{直角}{ちょっ|かく}より \R{小}{ちい}さいかな、 \R{大}{おお}きいかな」と \R{見}{み}て おくと、 まちがえません。
\R{上}{うえ}の \R{角}{かく}は \R{直角}{ちょっ|かく}より \R{小}{ちい}さいので、 90より \R{小}{ちい}さい 50°が \R{正}{ただ}しい（130°では ない）と わかります。

\begin{taisetsu}
\R{角}{かく}の \R{大}{おお}きさは、 まわった \R{大}{おお}きさ。\\
1\R{回転}{かい|てん}は 360°、 \R{直角}{ちょっ|かく}は 90°。
\end{taisetsu}

\bksubsection{5}{\R{時計}{と|けい}の みじかい はりの \R{角}{かく}}

\label{bk:1-20:tokei}\noindent\begin{minipage}[c]{0.36\linewidth}\centering
\begin{tikzpicture}[x=1mm, y=1mm]
  \bkclock{(0,0)}{14}{3}{0}
  \foreach \a/\c in {90/bkTheme,60/bkOrange,30/bkGreen} {\draw[\c, line width=1.6pt] (0,0) ++(\a:17) arc (\a:\a-30:17);}
  \node[font=\sffamily\bfseries\fontsize{8}{9}\selectfont, text=bkTheme] at (75:21) {30°};
  \node[font=\sffamily\bfseries\fontsize{8}{9}\selectfont, text=bkOrange!80!black] at (45:21.5) {30°};
  \node[font=\sffamily\bfseries\fontsize{8}{9}\selectfont, text=bkGreen!70!black] at (15:21.5) {30°};
\end{tikzpicture}
\end{minipage}\hfill
\begin{minipage}[c]{0.62\linewidth}
みじかい はりは、 12\R{時}{じ}から 3\R{時}{じ}までに \R{直角}{ちょっ|かく}（90°）だけ まわります。
その \R{間}{あいだ}に、 \R{数字}{すう|じ}の \R{間}{あいだ}が 3つ あります（12〜1、 1〜2、 2〜3）。
\R{数字}{すう|じ}の \R{間}{あいだ}は みんな \R{同}{おな}じ \R{大}{おお}きさなので、 90°を 3つに \R{分}{わ}けます。
30 ＋ 30 ＋ 30 ＝ 90 なので、 1つ\R{分}{ぶん}は 30°。
だから みじかい はりは、 1\R{時間}{じ|かん}に 30° まわります。
\end{minipage}

\par\vspace{3pt}
\noindent たしかめ：12\R{時}{じ}→3\R{時}{じ}→6\R{時}{じ}→9\R{時}{じ}→12\R{時}{じ}で \R{直角}{ちょっ|かく}が 4つ。 みじかい はりは 12\R{時間}{じ|かん}で ちょうど 1\R{回転}{かい|てん}（360°）します。

\Needspace{22\baselineskip}\bksubsection{6}{\R{分度器}{ぶん|ど|き}で \R{角}{かく}を かく}

\label{bk:1-20:kaku2}\R{分度器}{ぶん|ど|き}は、 \R{角}{かく}を かく ときにも つかえます。 60°の \R{角}{かく}を かいて みよう。

\noindent\begin{minipage}[c]{0.5\linewidth}\centering
\begin{tikzpicture}[x=1mm, y=1mm]
  \bkqsbundoki[5.4]{(0,0)}{28}{0,30,...,180}{0}
  \fill[bkRed] (60:26.5) circle (1.1);
  \draw[bkTheme, line width=1.3pt, dash pattern=on 3pt off 1.5pt] (0,0) -- (60:36);
  \node[anchor=south west, font=\sffamily\bfseries\fontsize{8}{9}\selectfont, text=bkRed!80!black, inner sep=0.5pt] at (62:29.5) {60の \R{点}{てん}};
\end{tikzpicture}
\end{minipage}\hfill
\begin{minipage}[c]{0.48\linewidth}
\begin{enumerate}[label={\bkcircled[bkTheme]{\arabic*}}, leftmargin=2\zw, itemsep=0pt, topsep=0pt]
\item まっすぐな \R{線}{せん}（\R{辺}{へん}）を 1\R{本}{ぽん} ひき、 はしを ちょう\R{点}{てん}に する。
\item \R{中心}{ちゅう|しん}を ちょう\R{点}{てん}に、 0の \R{線}{せん}を \R{辺}{へん}に \R{合}{あ}わせる。
\item その 0から かぞえて、 60の \R{目}{め}もりに \R{点}{てん}を うつ。
\item \R{分度器}{ぶん|ど|き}を はずし、 ちょう\R{点}{てん}と \R{点}{てん}を じょうぎで むすぶ。
\end{enumerate}
\end{minipage}

\par\vspace{3pt}
\noindent かけたら、 \R{分度器}{ぶん|ど|き}を もう\R{一度}{いち|ど} あてて、 60°に なって いるか たしかめよう。 \R{直角}{ちょっ|かく}より \R{小}{ちい}さい \R{角}{かく}に なって いれば、 \R{内}{うち}がわと \R{外}{そと}がわを まちがえて いません。

\par\vspace{6pt}
%% ---------------- p.206〜207 例題 ----------------
\begin{reidai}[\R{直角}{ちょっ|かく} いくつ\R{分}{ぶん}？]
\noindent\begin{minipage}[c]{0.5\linewidth}
\R{直角}{ちょっ|かく} 2つ\R{分}{ぶん}の \R{角}{かく}は \R{何度}{なん|ど}？\par\vspace{1mm}
\centering
\begin{tikzpicture}[x=1mm, y=1mm]
  \draw[bkInk!85, line width=1.3pt, line cap=round] (-16,0) -- (16,0) (0,0) -- (0,14);
  \bkqsright{(0,0)}{0}{2.6}\bkqsright{(0,0)}{90}{2.6}
  \draw[bkTheme, line width=1pt] (6,0) arc (0:180:6);
  \fill[bkInk] (0,0) circle (0.8);
\end{tikzpicture}
\end{minipage}\hfill
\begin{minipage}[c]{0.48\linewidth}
\kangaekata\\90° の 2つ\R{分}{ぶん}を たす。\\10が \mbox{9 ＋ 9 ＝ 18}こ。
\end{minipage}\par
\shiki 90 ＋ 90 ＝ 180 \kotaenext 180°
\end{reidai}

\begin{jibun}
\R{直角}{ちょっ|かく} 3つ\R{分}{ぶん}の \R{角}{かく}は \R{何度}{なん|ど}？\enspace{\small\color{bkOrange!58!black}ヒント：90は 10が 9こ}
\bkshikiran
\bkjans{90 ＋ 90 ＋ 90 ＝ 270 （10が \mbox{9 ＋ 9 ＋ 9 ＝ 27}こ）\quad 270°}
\end{jibun}

\begin{reidai}[\R{分度器}{ぶん|ど|き}で はかる]
\noindent\begin{minipage}[c]{0.5\linewidth}\centering
\begin{tikzpicture}[x=1mm, y=1mm]
  \bkqsbundoki{(0,0)}{27}{0,20,...,180}{0,40}
\end{tikzpicture}
\end{minipage}\hfill
\begin{minipage}[c]{0.48\linewidth}
この \R{角}{かく}は \R{何度}{なん|ど}？\par
\kangaekata\\\R{辺}{へん}は \R{右}{みぎ}の 0に \R{合}{あ}って いる。 だから \R{外}{そと}がわの \R{数}{かず}を 0、 10、 20、 … と よむ。
\end{minipage}\par
\noindent \R{直角}{ちょっ|かく}より \R{小}{ちい}さい \R{角}{かく}なので、 \R{答}{こた}えは 90より \R{小}{ちい}さい はず。\par
\kotae 40°\quad （\R{外}{そと}がわの \R{目}{め}もり。 140°では ない）
\end{reidai}

\begin{jibun}
\noindent\begin{minipage}[c]{0.5\linewidth}\centering
\begin{tikzpicture}[x=1mm, y=1mm]
  \bkqsbundoki{(0,0)}{24}{0,20,...,180}{0,70}
\end{tikzpicture}
\end{minipage}\hfill
\begin{minipage}[c]{0.48\linewidth}
この \R{角}{かく}は \R{何度}{なん|ど}？
\bkkotaeran[24mm]
\end{minipage}
\bkjans{70°}
\end{jibun}

\begin{reidai}[どちらの \R{目}{め}もり？]
\noindent\begin{minipage}[c]{0.5\linewidth}\centering
\begin{tikzpicture}[x=1mm, y=1mm]
  \bkqsbundoki{(0,0)}{25}{0,20,...,180}{180,60}
\end{tikzpicture}
\end{minipage}\hfill
\begin{minipage}[c]{0.48\linewidth}
この \R{角}{かく}は \R{何度}{なん|ど}？\par
\kangaekata\\\R{辺}{へん}は \R{左}{ひだり}の 0に \R{合}{あ}って いる。 だから \R{内}{うち}がわの \R{数}{かず}を よむ。
\end{minipage}\par
\noindent \R{直角}{ちょっ|かく}より \R{大}{おお}きい \R{角}{かく}なので、 \R{答}{こた}えは 90より \R{大}{おお}きい はず。\par
\kotae 120°\quad （\R{内}{うち}がわの \R{目}{め}もり。 60°では ない）
\end{reidai}

\begin{jibun}
\noindent\begin{minipage}[c]{0.5\linewidth}\centering
\begin{tikzpicture}[x=1mm, y=1mm]
  \node[font=\sffamily\bfseries] at (-24,22) {（1）};
  \bkqsbundoki{(0,0)}{22}{0,20,...,180}{180,30}
\end{tikzpicture}
\end{minipage}\hfill
\begin{minipage}[c]{0.48\linewidth}
この \R{角}{かく}は \R{何度}{なん|ど}？
\bkkotaeran[24mm]
\end{minipage}\par\vspace{4pt}
\noindent\begin{minipage}[c]{0.5\linewidth}\centering
\begin{tikzpicture}[x=1mm, y=1mm]
  \node[font=\sffamily\bfseries] at (-24,22) {（2）};
  \bkqsbundoki{(0,0)}{22}{0,20,...,180}{180,20}
\end{tikzpicture}
\end{minipage}\hfill
\begin{minipage}[c]{0.48\linewidth}
この \R{角}{かく}は \R{何度}{なん|ど}？
\bkkotaeran[24mm]
\end{minipage}
\bkjans{（1）150°\quad（2）160°（どちらも \R{左}{ひだり}の 0から \R{内}{うち}がわを よむ。 \R{直角}{ちょっ|かく}より \R{大}{おお}きい）}
\end{jibun}

\begin{jibun}
\R{数学}{すう|がく}の \R{道具}{どう|ぐ}⑥の じゅんで、 ノートに 40°と 120°の \R{角}{かく}を かこう。 かけたら \R{分度器}{ぶん|ど|き}で たしかめよう。
\bkjans{\R{中心}{ちゅう|しん}を ちょう\R{点}{てん}、 0の \R{線}{せん}を \R{辺}{へん}に \R{合}{あ}わせ、 その 0から 40（120）の \R{目}{め}もりに \R{点}{てん}を うって むすぶ。 40°は \R{直角}{ちょっ|かく}より \R{小}{ちい}さく、 120°は \R{大}{おお}きく なる}
\end{jibun}

\begin{reidai}[\R{時計}{と|けい}の はりが まわる \R{角}{かく}]
みじかい はりは 1\R{時間}{じ|かん}に 30°。 3\R{時間}{じ|かん}では \R{何度}{なん|ど}？\par
\noindent\begin{minipage}[c]{0.27\linewidth}\centering
\begin{tikzpicture}[x=1mm, y=1mm]
  \bkclock{(0,0)}{10}{3}{0}
  \draw[bkTheme, line width=1.2pt, -{Stealth[length=4pt]}] (0,0) ++(88:12.3) arc (88:2:12.3);
\end{tikzpicture}
\end{minipage}\hfill
\begin{minipage}[c]{0.72\linewidth}
\kangaekata 1\R{時間}{じ|かん}に 30°（\R{数字}{すう|じ} 1つ\R{分}{ぶん}）。 3\R{時間}{じ|かん}は 30° の 3つ\R{分}{ぶん}。\\
30は 10が 3つ。 10が \mbox{3 × 3 ＝ 9}こで 90。
\end{minipage}\par
\shiki 30 × 3 ＝ 90 \kotaenext 90°（\R{直角}{ちょっ|かく}！）
\end{reidai}

\begin{jibun}
みじかい はりは 1\R{時間}{じ|かん}に 30°。\par
（1）2\R{時間}{じ|かん}では \R{何度}{なん|ど}？
\bkshikiran
（2）5\R{時間}{じ|かん}では \R{何度}{なん|ど}？
\bkshikiran
\bkjans{（1）30 × 2 ＝ 60 （10が \mbox{3 × 2 ＝ 6}こ）\quad 60°\quad（2）30 × 5 ＝ 150 （10が \mbox{3 × 5 ＝ 15}こ）\quad 150°}
\end{jibun}

\begin{column}{どうして 1\R{回転}{かい|てん}は 360°？}
360は、 2でも 3でも 4でも 5でも 6でも、 ちょうど \R{分}{わ}けられる \R{数}{かず}です。 \R{角}{かく}を \R{分}{わ}けやすいように えらばれたと いわれて います。
\end{column}

\begin{column}{\R{長}{なが}い はりは？}
\R{長}{なが}い はりは、 60\R{分}{ぷん}で ぐるっと 1\R{回転}{かい|てん}（360°）します。
5\R{分}{ふん}で \R{数字}{すう|じ} 1つ\R{分}{ぶん}すすむので、 5\R{分}{ふん}で 30°。
15\R{分}{ふん}なら \R{数字}{すう|じ} 3つ\R{分}{ぶん}で 90°、 \R{直角}{ちょっ|かく}です。
\R{同}{おな}じ 30°でも、 みじかい はりは 1\R{時間}{じ|かん}、 \R{長}{なが}い はりは 5\R{分}{ふん}。 \R{長}{なが}い はりの ほうが ずっと はやく まわります。
\end{column}

\begin{reidai}[\R{直角}{ちょっ|かく}より \R{大}{おお}きい？ \R{小}{ちい}さい？]
\R{分度器}{ぶん|ど|き}で はかる \R{前}{まえ}に、 \R{直角}{ちょっ|かく}と くらべよう。 \R{直角}{ちょっ|かく}より \R{大}{おお}きいのは どれ？\par
\noindent\begin{minipage}[c]{0.56\linewidth}\centering
\begin{tikzpicture}[x=1mm, y=1mm]
  \foreach \i/\a/\t in {0/40/A,1/120/B,2/90/C,3/150/D} {
    \begin{scope}[shift={(\i*17+8,0)}]
      \bkqskaku{(0,0)}{0}{\a}{11}{4}
      \bkqstag{(0,15)}{\t}
    \end{scope}}
\end{tikzpicture}
\end{minipage}\hfill
\begin{minipage}[c]{0.42\linewidth}
\kangaekata \R{三角}{さん|かく}じょうぎの \R{直角}{ちょっ|かく}を あてる（1-19）。\\
すきまが あれば \R{大}{おお}きい。
\end{minipage}\par
\noindent A は じょうぎが はみ\R{出}{だ}す（\R{小}{ちい}さい）。 C は ぴったり（\R{直角}{ちょっ|かく}）。 B と D は すきまが ある（\R{大}{おお}きい）。\par
\kotae B と D
\end{reidai}

\begin{jibun}
\begin{bkcenter}
\begin{tikzpicture}[x=1mm, y=1mm]
  \foreach \i/\a/\t in {0/100/A,1/60/B,2/135/C,3/30/D} {
    \begin{scope}[shift={(\i*26+12,0)}]
      \bkqskaku{(0,0)}{0}{\a}{12}{4}
      \bkqstag{(-6,13)}{\t}
    \end{scope}}
\end{tikzpicture}
\end{bkcenter}
\noindent（1）\R{直角}{ちょっ|かく}より \R{小}{ちい}さいのは？\enspace\bkkakikomi[20mm]\hfill（2）\R{大}{おお}きいのは？\enspace\bkkakikomi[20mm]
\bkjans{（1）B と D\quad（2）A と C}
\end{jibun}

\par\vspace{6pt}
%% ---------------- p.208 まちがえやすいところ ----------------
\bksection{tsumazuki}{まちがえやすい ところ}

\begin{tsumazuki}{\R{辺}{へん}が \R{長}{なが}いと \R{角}{かく}も \R{大}{おお}きい？}
\begin{minipage}[c]{0.56\linewidth}
\R{角}{かく}の \R{大}{おお}きさは、 \R{辺}{へん}の \R{長}{なが}さでは きまりません。 ひらきぐあいで きまります。
\end{minipage}\hfill
\begin{minipage}[c]{0.42\linewidth}\centering
\begin{tikzpicture}[x=1mm, y=1mm]
  \bkqskaku{(0,0)}{0}{30}{26}{6}
  \node at (12,-4) {\bkbatsu};
  \node[font=\sffamily\footnotesize, text=bkInk] at (24,3) {A};
  \bkqskaku{(36,0)}{0}{60}{12}{5}
  \node at (42,-4) {\bkmaru};
  \node[font=\sffamily\footnotesize, text=bkInk] at (48,3) {B};
\end{tikzpicture}
\end{minipage}\par
\noindent{\small A より B の ほうが \R{大}{おお}きい \R{角}{かく}です。}
\end{tsumazuki}

\begin{tsumazuki}{\R{目}{め}もりの よみまちがい}
\begin{minipage}[c]{0.56\linewidth}
\R{辺}{へん}が \R{合}{あ}って いる 0から かぞえよう。\\
\bkbatsu\enspace 130°\hspace{2\zw}\bkmaru\enspace 50°
\end{minipage}\hfill
\begin{minipage}[c]{0.42\linewidth}\centering
\begin{tikzpicture}[x=1mm, y=1mm]
  \bkqsbundoki{(0,0)}{22}{0,50,90,180}{0,50}
\end{tikzpicture}
\end{minipage}
\end{tsumazuki}

\begin{yattemiyou}[おりがみで 45°を つくろう]
\noindent\begin{minipage}[c]{0.58\linewidth}
\begin{enumerate}[label={\bkcircled[bkGreen]{\arabic*}}, leftmargin=2.2\zw, itemsep=2pt, topsep=0pt]
\item おりがみの かどは \R{直角}{ちょっ|かく}（90°）。 \R{三角}{さん|かく}じょうぎで たしかめよう。
\item むかいあう かどと かどを \R{合}{あ}わせて、 ななめに \R{半分}{はん|ぶん}に おろう。
\item できた \R{三角形}{さん|かく|けい}の とがった かどを、 \R{分度器}{ぶん|ど|き}で はかろう。
\end{enumerate}
\end{minipage}\hfill
\begin{minipage}[c]{0.4\linewidth}\centering
\begin{tikzpicture}[x=1mm, y=1mm]
  \filldraw[fill=bkRoute6!30, draw=bkRoute6!70!black, line width=0.8pt] (0,0) rectangle (18,18);
  \draw[bkGray, line width=0.8pt, dash pattern=on 2pt off 1.5pt] (0,0) -- (18,18);
  \bkqsright{(0,0)}{0}{2.6}
  \draw[bk arrow, bkGray] (20,9) -- (26,9);
  \filldraw[fill=bkRoute6!30, draw=bkRoute6!70!black, line width=0.8pt, line join=round] (28,0) -- (46,0) -- (46,18) -- cycle;
  \draw[bkTheme, line width=1pt] (35,0) arc (0:45:7);
  \node[font=\sffamily\bfseries\fontsize{8}{9}\selectfont, text=bkTheme] at (39.5,-3) {？°};
\end{tikzpicture}
\end{minipage}
\par\vspace{4pt}
\noindent はかると 45°。 ななめに おると、 \R{直角}{ちょっ|かく}が \R{同}{おな}じ 2つに \R{分}{わ}かれます。 45 ＋ 45 ＝ 90 なので、 1つ\R{分}{ぶん}は 45°です。\par
\noindent \R{三角}{さん|かく}じょうぎの 45°の かどと ぴったり かさなるか、 あてて みよう。\enspace かさなった？\enspace\bkkakikomi[20mm]
\end{yattemiyou}

\begin{column}{\R{三角}{さん|かく}じょうぎの \R{角}{かく}を はかろう}
\begin{minipage}[c]{0.46\linewidth}
\R{分度器}{ぶん|ど|き}で はかると、\\\R{三角}{さん|かく}じょうぎの \R{角}{かく}は、\\90°・45°・60°・30°です。
\end{minipage}\hfill
\begin{minipage}[c]{0.52\linewidth}\centering
\begin{tikzpicture}[x=1mm, y=1mm]
  \bkqsjougi{(0,0)}{1}{22}{0}\bkqsright{(0,0)}{0}{3}
  \node[font=\sffamily\bfseries\scriptsize, text=bkTheme] at (16.5,2.2) {45°};
  \node[font=\sffamily\bfseries\scriptsize, text=bkTheme] at (2.8,16.5) {45°};
  \node[font=\sffamily\bfseries\scriptsize, text=bkTheme] at (5.6,5.6) {90°};
  \bkqsjougi{(28,0)}{2}{36}{0}\bkqsright{(28,0)}{0}{3}
  \node[font=\sffamily\bfseries\scriptsize, text=bkTheme] at (56.5,2.2) {30°};
  \node[font=\sffamily\bfseries\scriptsize, text=bkTheme] at (31,15.5) {60°};
  \node[font=\sffamily\bfseries\scriptsize, text=bkTheme] at (33.6,5.6) {90°};
\end{tikzpicture}
\end{minipage}
\end{column}

\clearpage
%% ---------------- p.209 つかわれているところ ----------------
\begin{butsuri}
\R{角}{かく}の \R{大}{おお}きさは、 きかいを \R{動}{うご}かす ときの 「ことば」です。
\begin{bkfigure}
\begin{tikzpicture}[x=1mm, y=1mm]
  % ロボットの うで
  \fill[white, rounded corners=3mm] (0,0) rectangle (39,40);
  \fill[bkGray!70, rounded corners=1mm] (5,8) rectangle (17,12);
  \draw[bkRoute6, line width=4pt, line cap=round] (11,13) -- (11,27);
  \draw[bkRoute6, line width=4pt, line cap=round] (11,27) -- (27,27);
  \fill[bkInk!75] (11,13) circle (2.2) (11,27) circle (2.2);
  \draw[bkInk!75, line width=1.3pt, line cap=round] (27,27) -- (31,30) (27,27) -- (31,24);
  \bkqsright{(11,27)}{270}{3.6}
  \node[font=\sffamily\bfseries\small, text=bkTheme] at (19,21.5) {90°};
  \node[font=\sffamily\bfseries\small] at (19.5,3.5) {ロボットの うで};
  % ソーラーパネル
  \begin{scope}[shift={(42,0)}]
  \fill[white, rounded corners=3mm] (0,0) rectangle (39,40);
  \draw[bkInk!60, line width=0.9pt] (3,10) -- (36,10);
  \fill[bkYellow] (32,33) circle (3.2);
  \foreach \a in {0,45,...,315} {\draw[bkYellow, line width=0.8pt] (32,33) ++(\a:4.2) -- ++(\a:1.8);}
  \draw[bkGray, line width=1.4pt] (24,10) -- (24,21);
  \draw[bkBlue!70!black, line width=3.2pt] (6,10) -- ++(30:27);
  \draw[bkTheme, line width=1pt] (16,10) arc (0:30:10);
  \node[font=\sffamily\bfseries\small, text=bkTheme] at (19.6,13.4) {30°};
  \node[font=\sffamily\bfseries\small] at (19.5,3.5) {ソーラーパネル};
  \end{scope}
  % スケートボード
  \begin{scope}[shift={(84,0)}]
  \fill[white, rounded corners=3mm] (0,0) rectangle (39,40);
  \bkkid{(19.5,15.5)}{17}{bkRoute2}{bkHair}
  \fill[bkStem, rounded corners=1.2mm] (9,13) rectangle (30,15.5);
  \fill[bkInk!80] (12.5,12) circle (1.4) (26.5,12) circle (1.4);
  \draw[bkTheme, line width=1.3pt, -{Stealth[length=5pt]}] (19.5,23) ++(200:14 and 9) arc (200:520:14 and 9);
  \node[font=\sffamily\bfseries\small, text=bkTheme] at (33,35) {360°};
  \node[font=\sffamily\bfseries\small] at (19.5,3.5) {スケートボード};
  \end{scope}
\end{tikzpicture}
\end{bkfigure}
\begin{itemize}[label={\color{bkPurple}●}, leftmargin=1.4\zw, itemsep=0pt, topsep=0pt]
\item ロボットの うでは、 「ひじを 90° まげる」の ように \R{角}{かく}で \R{動}{うご}かします。\\\R{手}{て}が どこに \R{来}{く}るかも、 \R{角}{かく}から \R{計算}{けい|さん}します。
\item ソーラーパネルは、 お\R{日}{ひ}さまの ほうへ \R{角}{かく}を きめて かたむけると、 \R{電気}{でん|き}が よく できます。
\item スケートボードの 「360」は、 1\R{回転}{かい|てん}（360°）する わざの \R{名前}{な|まえ}です。
\end{itemize}
\end{butsuri}

\bkfukidashi{delta}{まわる \R{角}{かく}は、 \R{第}{だい}4\R{巻}{かん}で 「\R{波}{なみ}」の \R{算数}{さん|すう}に つながるよ。\\「\R{回転}{かい|てん}と \R{波}{なみ}の \R{路線}{ろ|せん}」は まだまだ つづく！}

\begin{bkfigure}
\begin{tikzpicture}[x=1mm, y=1mm]
  \draw[bkRoute4, line width=3pt, line cap=round] (10,0) -- (122,0);
  \node[anchor=west, inner sep=0pt] at (0,12) {\bkrouteicon[1.4em]{4}\hspace{0.3em}{\sffamily\bfseries\small\color{bkRoute4!80!black}\R{回転}{かい|てん}と \R{波}{なみ}の \R{路線}{ろ|せん}}};
  \foreach \x/\c/\t/\s in {10/bkVol1/{1-15}/{なんしゅう},38/bkVol1/{1-19}/{\R{円}{えん}},66/bkVol1/{1-20}/{\R{角}{かく}},94/bkVol4/{4-8}/{まわる かげ},122/bkVol5/{5-15}/{\R{波}{なみ}の びぶん}} {
    \fill[white, draw=\c, line width=1.6pt] (\x,0) circle (2.6);
    \node[anchor=base, font=\sffamily\bfseries\footnotesize, text=\c!85!black] at (\x,-6.8) {\t};
    \node[anchor=base, font=\sffamily\footnotesize, text=bkInk] at (\x,-13.6) {\s};}
  \node[fill=bkTheme, text=white, rounded corners=1mm, font=\sffamily\bfseries\fontsize{8}{9}\selectfont, inner sep=1.5pt] at (66,5.5) {いま ここ};
\end{tikzpicture}
\end{bkfigure}

\clearpage
%% ---------------- p.210〜211 もんだい ----------------
\bkendproblems

\bkproblemgroup{きほん}{1}

\bkproblem \bkkakikomi に \R{数}{かず}を かきましょう。
\par\vspace{2pt}
\noindent\hspace*{1\zw}（1） \R{直角}{ちょっ|かく}は \bkkakikomi[14mm]°\hspace{1.5\zw}（2） \R{半回転}{はん|かい|てん}は \bkkakikomi[14mm]°\\
\hspace*{1\zw}（3） 1\R{回転}{かい|てん}は \bkkakikomi[14mm]°
\bkpans{（1）90°\quad（2）180°\quad（3）360°}

\bkproblem \R{角}{かく}の \R{大}{おお}きさを はかりましょう。
\begin{bkfigure}
\begin{tikzpicture}[x=1mm, y=1mm]
  \node[font=\sffamily\bfseries] at (-24,24) {（1）};
  \bkqsbundoki{(0,0)}{24}{0,20,...,180}{0,60}
  \begin{scope}[shift={(72,0)}]
  \node[font=\sffamily\bfseries] at (-24,24) {（2）};
  \bkqsbundoki{(0,0)}{24}{0,20,...,180}{180,40}
  \end{scope}
\end{tikzpicture}
\end{bkfigure}
\noindent\hspace*{1\zw}（1）\enspace\bkkakikomi[20mm]\hspace{6\zw}（2）\enspace\bkkakikomi[20mm]
\bkpans{（1）60°\quad（2）140°}

\bkproblem \R{直角}{ちょっ|かく} 4つ\R{分}{ぶん}の \R{角}{かく}は \R{何度}{なん|ど}ですか。\enspace{\small\color{bkOrange!58!black}ヒント：90は 10が 9こ}
\bkshikiran
\bkpans{90 ＋ 90 ＋ 90 ＋ 90 ＝ 360 （10が \mbox{9 × 4 ＝ 36}こ）\quad 360°}

\bkproblem \R{時計}{と|けい}の みじかい はりは、 1\R{時間}{じ|かん}に 30° まわります。\\4\R{時間}{じ|かん}では \R{何度}{なん|ど}？\enspace{\small\color{bkOrange!58!black}ヒント：30は 10が 3つ}
\bkshikiran
\bkpans{30 × 4 ＝ 120 （10が \mbox{3 × 4 ＝ 12}こ）\quad 120°}

\bkproblemgroup{チャレンジ}{1}

\bkproblem ロボットの うでを 90° まげて、 そこから また 90° まげました。 はじめから \R{何度}{なん|ど} まがりましたか。\enspace{\small\color{bkOrange!58!black}ヒント：90は 10が 9こ}
\bkshikiran
\bkpans{90 ＋ 90 ＝ 180 （10が \mbox{9 ＋ 9 ＝ 18}こ）\quad 180°}

\bkproblem \R{時計}{と|けい}の みじかい はりが、 12\R{時}{じ}から 6\R{時}{じ}まで まわる \R{角}{かく}は \R{何度}{なん|ど}？\enspace{\small\color{bkOrange!58!black}ヒント：30は 10が 3つ}
\bkshikiran
\bkpans{30 × 6 ＝ 180 （10が \mbox{3 × 6 ＝ 18}こ）\quad 180°}

\begin{furikaeri}
\begin{checklist}
\item \R{角}{かく}は 「まわった \R{大}{おお}きさ」と わかる。\hfill{\footnotesize\sffamily\color{bkTheme}→\pageref{bk:1-20:kaku}ページ}
\item 1\R{回転}{かい|てん}は 360°、 \R{直角}{ちょっ|かく}は 90°と わかる。\hfill{\footnotesize\sffamily\color{bkTheme}→\pageref{bk:1-20:do}ページ}
\item \R{分度器}{ぶん|ど|き}で \R{角}{かく}を はかれる。\hfill{\footnotesize\sffamily\color{bkTheme}→\pageref{bk:1-20:bundoki}ページ}
\end{checklist}
チェックが つかなかったら、 →の ページに もどろう。
\end{furikaeri}

\bkfukidashi{mina}{ノートの かどや、 ひらいた ドアの \R{角}{かく}を、 \R{分度器}{ぶん|ど|き}で はかって みよう！}
\clearpage
